\titledquestion{Multiple Choices}

Each question has \textbf{one or more} correct answer(s). Select all the correct answer(s). For each question, you will get 0 points if you select one or more wrong answers, but you will get 1 point if you select a non-empty subset of the correct answers.

Write your answers in the following table.

\begin{table}[h]
    \centering
    \renewcommand{\arraystretch}{1.25}
    \begin{tabular}{|p{2cm}|p{2cm}|p{2cm}|p{2cm}|p{2cm}|}
        \hline 
        (a) & (b) & (c) & (d) & (e)\\
        \hline
        %%%%%%%%%%%%%%%%%%%%%%%%%%%%%%%%%%%%%%%%%%%%%%%%%%%%%%%%%%
        % YOUR ANSWER HERE.
        ACD & ABC & AC & AC & BC \\
        %%%%%%%%%%%%%%%%%%%%%%%%%%%%%%%%%%%%%%%%%%%%%%%%%%%%%%%%%%
        \hline
    \end{tabular}
\end{table}

\begin{parts}

%=========================================================
% (a)
%=========================================================

\part[3] Which of the following sorting algorithm(s) is(are) stable?

\begin{choices}
    \choice Bubble-Sort

    \choice Quick-Sort (picking the first element as pivot)

    \choice Insertion-Sort

    \choice Merge-Sort
\end{choices}


%=========================================================
% (b)
%=========================================================

\part[3]
Given $n$ distinct elements
\[
a_1<a_2<\cdots<a_n,
\]
suppose Randomized Quick-Sort chooses the pivot uniformly at random from
the current subarray. Which of the following statement(s) is(are)
\textbf{true}?

\begin{choices}

    \choice
    Any two adjacent elements $a_i$ and $a_{i+1}$ will be compared during
    the execution of Randomized Quick-Sort.

    \choice
    For $i<j$, the probability that $a_i$ and $a_j$ are compared is
    \[
    \frac{2}{j-i+1}.
    \]

    \choice
    The probability that $a_1$ and $a_n$ are compared is
    \[
    \frac{2}{n}.
    \]

    \choice
    The expected number of comparisons made by Randomized Quick-Sort is
    $\Theta(n^2)$.

\end{choices}

%=========================================================
% (c)
%=========================================================

\part[3]
Which of the following statement(s) about Quick-Sort is(are)
\textbf{true}?

\begin{choices}

    \choice
    The median-of-three strategy chooses the median of the first, middle,
    and last entries as the pivot. It may give a better approximation of
    the true median, but it does not guarantee a balanced partition.

    \choice
    If Quick-Sort uses in-place partitioning, then it uses
    $\Theta(1)$ extra memory in every case.

    \choice
    For sufficiently small subarrays, using Insertion-Sort instead of
    continuing the Quick-Sort recursion can be a reasonable optimization.

    \choice
    The recursion depth of Quick-Sort is always $\Theta(\log n)$,
    regardless of how the pivots are chosen.

\end{choices}

%=========================================================
% (d)
%=========================================================

\part[3]
Which of the following statement(s) about divide-and-conquer algorithms
and their applications is(are) \textbf{true}?

\begin{choices}

    \choice
    A divide-and-conquer algorithm usually consists of three steps:
    dividing the original problem into smaller subproblems, recursively
    solving the subproblems, and combining their solutions.

    \choice
    Any recursive algorithm can be considered a divide-and-conquer algorithm
    because it always solves smaller instances of the original problem.

    \choice
    FFT applies the divide-and-conquer idea by recursively decomposing a
    polynomial into its even-degree and odd-degree parts.

    \choice
    If a divide-and-conquer algorithm divides a problem of size $n$
    into two subproblems of size $n/2$, its running time must be
    $\Theta(n\log n)$, regardless of the cost of combining the solutions.

\end{choices}

%=========================================================
% (e)
%=========================================================



\part[3]
Which of the following statement(s) about divide-and-conquer
sorting algorithms is(are) \textbf{true}?

\begin{choices}

    \choice
    In merge sort, if every element in the left half is smaller
    than every element in the right half, the recursive sorting
    of both halves can be skipped.
    
    \choice
    After a correct quicksort partition, every element in the left
    partition is no greater than any element in the right partition,
    even though neither partition is necessarily sorted.


    \choice
    When merging two sorted subarrays of lengths $p$ and $q$,
    the standard merge procedure requires at most
    $p+q-1$ element comparisons.

    \choice
    If quicksort consistently divides a problem of size $n$
    into subproblems of approximately $n/4$ and $3n/4$,
    its running time is $\Theta(n^2)$ because the partitions
    are unbalanced.

\end{choices}

\end{parts}